%% ---- 1. Introdução ---------------------------------------------------------
\section{Introdução}

Um recomendador sequencial recebe o histórico de um usuário e devolve uma lista
de itens. Para um item $v$ dessa lista, uma pergunta natural é \emph{quais
itens do histórico explicam a presença de $v$}. Uma resposta contrafactual é um
conjunto de posições do histórico com duas propriedades: se esses itens forem
trocados por itens neutros, $v$ sai da lista (o conjunto é \emph{necessário});
se apenas eles forem mantidos, $v$ continua na lista (o conjunto é
\emph{suficiente}). Queremos também que o conjunto seja pequeno, porque uma
explicação com metade do histórico não explica nada.

Esses três critérios não são comensuráveis, e o maior conjunto possível (o
histórico inteiro) é ótimo nos dois primeiros e péssimo no terceiro. Os métodos
de caixa-preta existentes para recomendação sequencial combinam os critérios em
um único número e devolvem uma única
resposta~\cite{scarcelli2025gece,kim2026trace}; quem escolhe o compromisso é o
peso da soma, fixado antes da busca.

Neste trabalho tratamos o problema como multiobjetivo. A busca é sobre uma
máscara binária das posições do histórico, o modelo é consultado apenas por
seus escores, e a saída é uma frente de Pareto de (necessidade, suficiência,
tamanho), aproximada por algoritmos evolutivos multiobjetivo, a começar pelo
NSGA-II~\cite{deb2002nsga2}. Os recomendadores explicados são dois
SASRec~\cite{kang2018sasrec}, treinados no MovieLens-20M e no Zvuk (música);
eles são o objeto da explicação, não a contribuição, e seu treino está no
Apêndice~\ref{ap:treino}.

As contribuições são:
\begin{itemize}
  \item uma formulação com necessidade e suficiência separadas, estimadas por
        intervenções repetidas com fundos fixos, o que torna a aptidão
        determinística (Seções~\ref{sec:formulacao} e~\ref{sec:mc});
  \item uma comparação sob o mesmo orçamento de consultas ao modelo entre
        quatro algoritmos multiobjetivo (NSGA-II, NSGA-III, SMS-EMOA e MOEA/D),
        um algoritmo genético de objetivo único com os mesmos operadores, um
        guloso, valores de Shapley, quatro heurísticas, a enumeração exata de
        conjuntos pequenos e o GECE (Seção~\ref{sec:resultados});
  \item estudos de operadores, de uma população inicial semeada, do valor do
        objetivo de suficiência e da distribuição de fundo, uma verificação
        contra a frente exata e a replicação em um segundo conjunto de
        instâncias e em um segundo conjunto de dados.
\end{itemize}

O resultado tem três partes. Entre explicações de até 10 posições, os
algoritmos baseados em dominância de Pareto (NSGA-II, NSGA-III e SMS-EMOA) são
os melhores métodos; o MOEA/D, que decompõe o problema em subproblemas escalares, e o
algoritmo genético escalarizado ficam muito atrás. Quando máscaras de qualquer
tamanho contam, métodos que percorrem todos os tamanhos (guloso, Shapley)
superam o NSGA-II, e semear a população inicial com um ranking por ablação
inverte esse resultado ao custo de uma perda pequena nas explicações pequenas.
Necessidade e suficiência conflitam nas explicações pequenas, nos dois
conjuntos de dados, e esse conflito não é ruído de estimação. A comparação
entre métodos se repete em 99 instâncias novas e no Zvuk.
